% GENERATED FILE. Edit canonical lessons, then run npm run generate:books.
% canonical-source-hash: fnv1a64:c135067668777cc1
% canonical-lessons: RU-C270-ya-by-khotel, RU-C270-yesli-by, RU-C270-review, RU-W270-pismo

\chapter{Would, and a Letter About a Trip}
\label{ch:ru-would-and-a-letter-about-a-trip}
\hlchaptermodality{270}

\begin{chapteropening}
\textbf{By the end of this chapter:} \emph{I can say what I would like with \ru{я} \ru{бы} \ru{хотел}, what would happen with \ru{если} \ru{бы}, and write six connected lines about a trip that end on a condition.}

Complete RU-W270-pismo: six lines to one reader about a trip, the scene in the imperfective, the events in order in the perfective, a feeling, and a last line with \ru{бы}.
\end{chapteropening}

\section[\texorpdfstring{\ru{я} \ru{бы} \ru{хотел}, \ru{я} \ru{бы} … (ya by khotél, ya …)}{я бы хотел, я бы … (ya by khotél, ya …)}]{\ru{я} \ru{бы} \ru{хотел} — would, and the polite request}

\label{lesson:RU-C270-ya-by-khotel}



\begin{quote}
Say \textquotedblleft{}I want a coffee\textquotedblright{}. (\emph{Ya khochú kófe.}) Then \textquotedblleft{}if I have time, I'll ring\textquotedblright{}. (\emph{Yésli búdet vrémya, ya pozvonyú.}) That was a real plan. Now a softer world: what would be.
\end{quote}

\begin{grammarlens}[title={Grammar Lens: \ru{бы} and the -\ru{л} form}]
\textbf{\ru{бы}} — \emph{by} — is a small word that makes \textquotedblleft{}would\textquotedblright{}. It goes with the past form of the verb, the \emph{-\ru{л}} form, and here that form says nothing about the past:

\begin{quote}
\textbf{\ru{Я} \ru{бы} \ru{хотел} \ru{кофе}.} — \emph{ya by khotél kófe} — I'd like a coffee. \textbf{\ru{Я} \ru{бы} \ru{хотела} \ru{чай}.} — \emph{ya by khotéla chay} — I'd like tea; a woman says it. \textbf{\ru{Я} \ru{бы} \ru{пошёл} \ru{в} \ru{кино}.} — \emph{ya by poshól v kinó} — I'd go to the cinema.
\end{quote}

\textbf{\ru{Бы}} can also follow the verb: \emph{\ru{Я} \ru{хотел} \ru{бы} \ru{кофе}.} It agrees with nothing; the \emph{-\ru{л}} form does the agreeing.
\end{grammarlens}

\begin{grammarlens}[title={Grammar Lens: the polite request}]
\emph{\ru{Я} \ru{хочу}} is direct; \textbf{\ru{я} \ru{бы} \ru{хотел}} is the polite way to ask for a thing. For asking someone to do something, Russian says \textquotedblleft{}could you not …\textquotedblright{}:

\begin{quote}
\textbf{\ru{Вы} \ru{не} \ru{могли} \ru{бы} \ru{говорить} \ru{медленно}?} — \emph{vy ne moglí by govorít médlenno?} — Could you speak slowly? With \emph{\ru{повторить}}: \textbf{\ru{Вы} \ru{не} \ru{могли} \ru{бы} \ru{повторить}?}
\end{quote}

\textbf{\ru{Могли}} is the \emph{-\ru{ли}} form of \emph{\ru{мочь}}, for polite \emph{\ru{вы}}; \emph{\ru{повторить}} is the partner of \emph{\ru{повторять}}. The \textbf{\ru{не}} makes it gentler, not negative.
\end{grammarlens}

\begin{cousinweb}
\textbf{\ru{бы}} is, by the usual account, a worn-down old form of \textbf{\ru{быть}}, left over from a past tense Russian no longer has; EtymDB 2.1 lists \emph{\ru{бы}} without a parent. It is the same \emph{\ru{бы}} that sits inside \textbf{\ru{чтобы}}, \textquotedblleft{}in order to\textquotedblright{}, a wish fastened to \emph{\ru{что}}.
\end{cousinweb}

\subsection*{Your turn}
\begin{itemize}
\raggedright
  \item \emph{Say it:} \emph{Ya by khotél kófe}, then a woman: \emph{Ya by khotéla kófe.}
  \item \emph{Say it:} \emph{Ya khotél by chay.} The \emph{by} after the verb
  \item \emph{Say it:} \emph{Vy ne moglí by govorít médlenno?}
  \item \emph{Answer:} at a café the waiter asks \emph{Shto vy khotítye?} Answer politely with \emph{ya by khotél} or \emph{ya by khotéla}
\end{itemize}

\begin{tcolorbox}[breakable,colback=teal!4,colframe=teal!35!black,title={Before you move on}]
\textquotedblleft{}I'd like a coffee\textquotedblright{}, a woman? (\emph{\textbf{\ru{Я} \ru{бы} \ru{хотела} \ru{кофе}.}}) Which verb form goes with \emph{\ru{бы}}? (\textbf{The -\ru{л} form}, but it is not past.)
\end{tcolorbox}
\section[\texorpdfstring{\ru{если} \ru{бы} \ru{у} \ru{меня} … (yésli by u menyá …)}{если бы у меня … (yésli by u menyá …)}]{\ru{если} \ru{бы} — the condition that is not so}

\label{lesson:RU-C270-yesli-by}



\begin{quote}
Say \textquotedblleft{}I'd like a coffee\textquotedblright{}. (\emph{Ya by khotél kófe.}) Then the real plan: \textquotedblleft{}if I have time, I'll ring\textquotedblright{}. (\emph{Yésli búdet vrémya, ya pozvonyú.})
\end{quote}

\begin{grammarlens}[title={Grammar Lens: \ru{у} \ru{меня} \ru{было}, I had}]
You have said \textbf{\ru{у} \ru{меня} \ru{проблема}}, \textquotedblleft{}at me, a problem\textquotedblright{}. \textquotedblleft{}I have time\textquotedblright{} is built the same way: \textbf{\ru{у} \ru{меня} \ru{есть} \ru{время}}. In the past, the thing you had is the subject, and \textbf{\ru{было}} agrees with it: \emph{\ru{время}} is neuter, so \textbf{\ru{у} \ru{меня} \ru{было} \ru{время}} — \emph{u menyá býlo vrémya} — I had time.
\end{grammarlens}

\begin{grammarlens}[title={Grammar Lens: \ru{бы} on both sides}]
\begin{quote}
\textbf{\ru{Если} \ru{бы} \ru{у} \ru{меня} \ru{было} \ru{время}, \ru{я} \ru{бы} \ru{позвонил}.} — \emph{yésli by u menyá býlo vrémya, ya by pozvoníl} — If I had time, I would ring. \textbf{\ru{Если} \ru{бы} \ru{я} \ru{знал}, \ru{я} \ru{бы} \ru{помог}.} — \emph{yésli by ya znal, ya by pomóg} — If I had known, I would have helped.
\end{quote}

\textbf{\ru{Бы}} stands on both sides, each time with the past form: \emph{\ru{знал}}, \emph{\ru{позвонил}}, and \emph{\ru{помог}}, one of the few without an \emph{\ru{л}}. Russian marks no time here: the first can mean \textquotedblleft{}if I had time now\textquotedblright{} or \textquotedblleft{}if I had had time then\textquotedblright{}, and the story around it decides. \textbf{\ru{Помог}}, \textquotedblleft{}helped\textquotedblright{}, is the past of \emph{\ru{помочь}}, the partner of \emph{\ru{помогать}}.
\end{grammarlens}

\begin{culture}
\begin{quote}
\textbf{\ru{Если} \ru{будет} \ru{время}, \ru{я} \ru{позвоню}.} — It may well happen: a plan. \textbf{\ru{Если} \ru{бы} \ru{было} \ru{время}, \ru{я} \ru{бы} \ru{позвонил}.} — It is not so: a wish, or a regret.
\end{quote}
\end{culture}

\subsection*{Your turn}
Say these aloud:

\begin{itemize}
\raggedright
  \item \emph{U menyá býlo vrémya.}
  \item \emph{Yésli by u menyá býlo vrémya, ya by pozvoníl.}
  \item \emph{Yésli by ya znal, ya by pomóg}, then a woman: \emph{Yésli by ya znála, ya by pomoglá.}
  \item one wish of your own: \emph{Yésli by u menyá býlo vrémya, ya by …}
\end{itemize}

\begin{tcolorbox}[breakable,colback=teal!4,colframe=teal!35!black,title={Before you move on}]
\textquotedblleft{}I had time\textquotedblright{}? (\emph{\textbf{\ru{У} \ru{меня} \ru{было} \ru{время}.}}) A plan or a wish: \emph{\ru{если} \ru{бы} \ru{было} \ru{время}}? (\textbf{A wish}: it is not so.) How many times does \emph{\ru{бы}} appear? (\textbf{Twice}, once on each side.)
\end{tcolorbox}
\section[(review)]{Review — what will be, and what would be}

\label{lesson:RU-C270-review}



\begin{quote}
Ask politely for a tea, as a woman. (\emph{Ya by khotéla chay.}) Then ask someone to speak slowly. (\emph{Vy ne moglí by govorít médlenno?})
\end{quote}

\begin{grammarlens}[title={What you've built}]
Two lessons, and Russian can now say what is not so. \textbf{\ru{Бы}} with the \emph{-\ru{л}} form, whatever the time. \textbf{\ru{Я} \ru{бы} \ru{хотел}} to ask for a thing, \textbf{\ru{вы} \ru{не} \ru{могли} \ru{бы}} to ask for a favour. And two kinds of \emph{\ru{если}}: the future for a plan, \textbf{\ru{бы}} on both sides for a wish.

The same dream, two ways:

\begin{quote}
\textbf{\ru{Если} \ru{у} \ru{меня} \ru{будут} \ru{деньги}, \ru{я} \ru{буду} \ru{путешествовать}.} — a plan. \textbf{\ru{Если} \ru{бы} \ru{у} \ru{меня} \ru{были} \ru{деньги}, \ru{я} \ru{бы} \ru{путешествовал}.} — a wish.
\end{quote}

\emph{\ru{Деньги}}, \textquotedblleft{}money\textquotedblright{}, is plural, so \emph{\ru{будут}} and \emph{\ru{были}}. \emph{\ru{Путешествовать}} is imperfective, so its future is \emph{\ru{буду}} and the infinitive.
\end{grammarlens}

\subsection*{Your turn}
\begin{itemize}
\raggedright
  \item \emph{Say it:} \emph{Yésli by u menyá býli déngi, ya by puteshéstvoval.}
  \item \emph{Say it:} \emph{Yésli by ya znála, ya by pomoglá.}
  \item \emph{Answer:} \emph{Ty búdesh puteshéstvovat létom?} \emph{Éto zavísit ot déneg}, then one line with \emph{yésli}
  \item \emph{Say it:} \emph{Ya mechtáyu o móre. Yésli by u menyá býlo vrémya, ya by puteshéstvoval.}
\end{itemize}

\begin{tcolorbox}[breakable,colback=teal!4,colframe=teal!35!black,title={Before you move on}]
A plan, or a wish: \emph{\ru{если} \ru{бы} \ru{у} \ru{меня} \ru{были} \ru{деньги}}? (\textbf{A wish.}) \textquotedblleft{}Could you …?\textquotedblright{}, politely? (\emph{\textbf{\ru{Вы} \ru{не} \ru{могли} \ru{бы} …?}}) \textquotedblleft{}I had time\textquotedblright{}? (\emph{\textbf{\ru{У} \ru{меня} \ru{было} \ru{время}.}})
\end{tcolorbox}
\section[(six lines)]{Six lines about a trip}

\label{lesson:RU-W270-pismo}



\begin{quote}
The four-line writing step held together with \emph{\ru{и}}, \emph{\ru{но}}, \emph{\ru{а}} and \emph{\ru{потому} \ru{что}}. This time there are six, and they tell a story, as Anna's rainy Saturday did.

Pick one person who will read it: \emph{\ru{друг}}, \emph{\ru{подруга}}, someone in your family.
\end{quote}

\subsection*{What to know first}
A told trip is more than a list of events. It has a \textbf{scene}, in the imperfective: \emph{\ru{было}}, \emph{\ru{шёл} \ru{дождь}}, \emph{\ru{мы} \ru{жили} \ru{в} \ru{гостинице}}. It has \textbf{steps} in order, each one finished, in the perfective. It has a \textbf{feeling}. And a good one ends by looking forward, or wishing.

A past verb agrees with who did it: \emph{-\ru{л}}, \emph{-\ru{ла}}, \emph{-\ru{ли}}, or, in a few such as \emph{\ru{мог}}, the same endings without the \emph{\ru{л}}: \emph{\ru{мог}}, \emph{\ru{могла}}, \emph{\ru{могли}}. A verb after \emph{\ru{бы}} takes the same forms.

\subsection*{Writing — connected composition}
On paper, write \textbf{six lines} to your reader about a trip, real or invented, one line for each of these:

\begin{enumerate}
\raggedright
  \item the scene: when, with \emph{\ru{назад}} (\emph{\ru{год} \ru{назад}}), and what it was like, in the imperfective;
  \item the first step, with \emph{\ru{сначала}}, in the perfective;
  \item the next step, with \emph{\ru{потом}} or \emph{\ru{затем}}, and a reason, with \emph{\ru{потому} \ru{что}}, \emph{\ru{поэтому}} or \emph{\ru{так} \ru{что}}; a problem works well here: \emph{\ru{я} \ru{опоздал} \ru{на} \ru{поезд}};
  \item something that cut in, with \emph{\ru{когда}} or \emph{\ru{пока}}, and how it ended, with \emph{\ru{наконец}};
  \item how you felt, with \emph{\ru{я} \ru{был} \ru{рад}, \ru{что} …}, \emph{\ru{я} \ru{устал}} or \emph{\ru{мне} \ru{было} …}, or what you thought, with \emph{\ru{по}-\ru{моему}} and a comparison: \emph{\ru{лучше}, \ru{чем} …};
  \item one line forward, with \emph{\ru{если} \ru{бы} \ru{у} \ru{меня} \ru{было} \ru{время}, \ru{я} \ru{бы} …}, or \emph{\ru{я} \ru{бы} \ru{хотел} …}, or a hope: \emph{\ru{я} \ru{надеюсь}, \ru{что} …}, \emph{\ru{я} \ru{мечтаю} …}.
\end{enumerate}

Use only words you already write. The exercise is the joining, and the aspect of each verb.

\begin{tcolorbox}[breakable,colback=teal!4,colframe=teal!35!black,title={Before you move on}]
Now go through your lines and find every verb. In line one, is it imperfective, a scene? In lines two, three and four, is each step perfective, one deed done? That is the slip to check first.

Then the endings: \emph{-\ru{л}} for a man, \emph{-\ru{ла}} for a woman, \emph{-\ru{ли}} for more than one, after \emph{\ru{бы}} as well.

Last, the commas: one before \emph{\ru{что}}, \emph{\ru{потому} \ru{что}} and \emph{\ru{чем}}, one before \emph{\ru{поэтому}} or \emph{\ru{так} \ru{что}} when they join inside a sentence, and one at the seam after a \emph{\ru{когда}}, \emph{\ru{пока}} or \emph{\ru{если}} half. If line six uses \emph{\ru{если} \ru{бы}}, does \emph{\ru{бы}} stand on both sides of the comma?
\end{tcolorbox}
